\documentclass[11pt,a4paper]{article}
\usepackage[T1]{fontenc}
\usepackage[utf8]{inputenc}
\usepackage{amsmath,amssymb,amsthm}
\usepackage[margin=2.7cm]{geometry}
\usepackage{booktabs}

\usepackage[expansion=false]{microtype}
\usepackage{xcolor}
\usepackage[colorlinks=true,linkcolor=blue!50!black,citecolor=blue!50!black,urlcolor=blue!50!black]{hyperref}

\newtheorem{theorem}{Theorem}[section]
\newtheorem{proposition}[theorem]{Proposition}
\newtheorem{lemma}[theorem]{Lemma}
\newtheorem{corollary}[theorem]{Corollary}
\newtheorem{conjecture}[theorem]{Conjecture}
\newtheorem{hypothesis}[theorem]{Hypothesis}
\theoremstyle{definition}
\newtheorem{definition}[theorem]{Definition}
\newtheorem{algorithm}[theorem]{Algorithm}
\theoremstyle{remark}
\newtheorem{remark}[theorem]{Remark}
\newtheorem{observation}[theorem]{Observation}

\newcommand{\slog}{\operatorname{slog}}
\newcommand{\R}{\mathbb{R}}
\newcommand{\Z}{\mathbb{Z}}
\newcommand{\N}{\mathbb{N}}
\newcommand{\Sone}{S^1}
\newcommand{\Ph}[1]{\widehat{P}(#1)}
\newcommand{\eff}{\mathrm{eff}}

\title{Fine Structure of Mixed-Base Tetration Phases:\\
Certified Evaluation, a Universal Mode-Decay Law,\\
and a Boundary Essential Singularity\\[2mm]
\large Paper IV of the mixed-base tetration series}
\author{Janis Justus\thanks{See the Authorship and AI Contribution Statement at the end of this
manuscript: the research reported here was carried out in a directed human--AI mode, with the
author acting primarily as research director and the mathematical development and numerical
campaign largely produced by AI systems under his direction.}}
\date{July 2026}

\begin{document}
\maketitle

\begin{abstract}
For regular tetration branches $T_b$, $T_d$ with $b,d>e^{1/e}$, the mixed-base phase
$\Phi_{b,d}(\theta)=\lim_{n\to\infty}\bigl(\slog_b(T_d(n+\theta))-(n+\theta)\bigr)$ is a small
$1$-periodic circle map whose existence, groupoid structure and $C^1$-diffeomorphism property
were established in the companion papers \cite{phaselaw,unified}. This paper resolves, partly
by proof and partly by a high-precision measurement campaign, the fine structure of these
phases, and turns them into a certified computational primitive for base change in tetration.

On the theoretical side we prove: an exact two-level entry and an a-posteriori error
certificate for the logarithmic peel ladder, making every phase evaluation carry a proven
error bound; a regularity dictionary ``branch $C^m$ $\Rightarrow$ Fourier modes
$\asymp k^{-(m+2)}$'' with an explicit seam-jump formula, calibrated on explicit $C^1$
branches (measured exponents $3.06$ and $3.006$); a first-order perturbation theory in the
base mismatch $\varepsilon=\ln\log_b d$, in which the derivative of the phase at the diagonal
$b=d$ is given by a response kernel solving a linear cohomological equation, explaining the
empirical amplitude law $A_k(b,d)\approx C_k\,\lvert\ln\log_b d\rvert$; an exact hub
decomposition $\mu_{b,d}=m(d)-m(b)+\Delta$ of the mean shifts with second- and third-order
defect formulas validated to $10^{-15}$; and a parabolic passage-time derivation of the edge
law $m(b)\sim-\pi\sqrt{2\eta/e}\,(b-\eta)^{-1/2}$ as $b\downarrow\eta=e^{1/e}$, matching the
measured coefficient to $2.3\%$.

On the empirical side, Fourier modes of the regular phases were measured to $k=300$ after
identifying and quantifying an instrument systematic of the underlying Kneser engine
(usable band $\approx 2.3\,n_{\mathrm{lim}}$ theta harmonics). The decay profile
$\sigma_\eff(k)$ is universal across all fourteen measured base pairs to three to four
decimals; a direct exploration of the complexified phase locates an isolated
\emph{essential singularity} per period on the strip boundary, at pair-universal height
$y_*\approx 0.055$ and with pair-universal strength, only its horizontal position $x_*(b,d)$
depending on the pair. Steepest-descent asymptotics link the singularity data to the measured
four-parameter decay law and yield mode-count formulas for base-change tables
($\approx 300$ modes for $150$ digits, $\approx 1400$ for $500$ digits). We close with a
hybrid data format---mean-shift hub, Fourier head plus universal curve, certified peel ladder
beyond---in which cross-base identities serve as proof residuals rather than heuristic checks.
\end{abstract}

\section{Introduction}\label{sec:intro}

Regular tetration to a base $a>e^{1/e}$ is the increasing real-analytic solution $T_a$ of
\begin{equation}\label{eq:tetration}
T_a(0)=1,\qquad T_a(x+1)=a^{\,T_a(x)},
\end{equation}
in the Kneser--Paulsen--Cowgill normalization \cite{kneser,kouznetsov,paulsencowgill,paulsen},
with Abel inverse $A_a=T_a^{-1}=\slog_a$. Given two admissible bases, the mixed-base height map
$F_{b,d}=A_b\circ T_d$ converts $d$-height into $b$-height. The first paper of this series
\cite{phaselaw} proved, conditionally on explicit channel and peeling-growth estimates for the
regular branch, the \emph{phase law}
\begin{equation}\label{eq:phaselaw}
\slog_b\bigl(T_d(x)\bigr)-x \;=\; \Phi_{b,d}(\{x\})+o(1),\qquad x\to+\infty,
\end{equation}
with $\Phi_{b,d}$ continuous and $1$-periodic, decomposed as
$\Phi_{b,d}=\mu_{b,d}+P_{b,d}$ into a mean shift and a centered periodic profile, and
established the exact groupoid identity $F_{b,d}=F_{b,c}\circ F_{c,d}$ together with its
nonlinear phase cocycle. The companion note \cite{unified} then de-conditionalized the
smoothness assumptions: the lifted phase maps $\widetilde H_{b,d}=\mathrm{id}+\Phi_{b,d}$ form a
groupoid of orientation-preserving circle homeomorphisms with
$\widetilde H_{b,d}^{-1}=\widetilde H_{d,b}$, the inverse mean antisymmetry
$\mu_{b,d}=-\mu_{d,b}$ holds exactly with no smoothness at all, and under a $C^1$ reading of
the channel hypotheses the phase maps are $C^1$ circle diffeomorphisms with the exact
derivative product identity
\begin{equation}\label{eq:product}
\widetilde H_{b,d}'(\theta)\;=\;\frac{\alpha\,T_d'(\theta)}{T_b'\bigl(\widetilde
H_{b,d}(\theta)\bigr)}\;\prod_{j=1}^{\infty}\frac{\alpha\,T_d(j+\theta)}{T_b\bigl(j+\widetilde
H_{b,d}(\theta)\bigr)}\;>\;0,\qquad \alpha=\log_b d .
\end{equation}

Three questions were left open, and they are precisely the questions that decide whether the
phases can be used as a stable computational primitive for base change:

\begin{itemize}
\item[(Q1)] \emph{Regularity and mode decay.} Is $\Phi_{b,d}$ real-analytic on a strip
(\cite[Conj.~8.1]{phaselaw}), merely $C^\infty$, or something in between? Equivalently: how do
the Fourier amplitudes $A_k=2\lvert\widehat P_{b,d}(k)\rvert$ decay? The answer fixes the size
of any tabulated representation of $\Phi$.
\item[(Q2)] \emph{Laws for the shift and the amplitudes.} The data of \cite{phaselaw} hinted at
$A_1\approx C_1\lvert\ln\alpha\rvert$ with a near-universal constant $C_1$, and at near-additive
mean shifts. Are there provable first-order laws behind these observations, and what is the
exact bookkeeping for $\mu_{b,d}$ across many bases?
\item[(Q3)] \emph{Certified evaluation.} Can every evaluation of $\Phi$, of its inverse, and of
every reconstructed pair quantity carry a proven error bound suitable for use inside an
automated verification loop?
\end{itemize}

This paper answers (Q3) affirmatively by proof, answers (Q2) by a perturbation theory whose
predictions are confirmed to many digits, and answers (Q1) by a measurement campaign that
replaces the dichotomy ``strip versus $C^\infty$'' with a sharper picture: the phase extends
holomorphically to a strip \emph{punctured} at one point per period, where it has an isolated
essential singularity of pair-universal type, whose height $y_*\approx0.055$ (combined
estimate $\sigma_*=0.058\pm0.006$) is the strip
half-width. All numerical claims come from a July 2026 measurement campaign against the
author's high-precision Kneser engine (a fork of Levenstein's \texttt{fatou.gp}
\cite{fatou}); Appendix~\ref{app:repro} documents the pipeline and the proof-residual gates
under which the data were accepted.

\subsection{Contributions}

\emph{Certified algorithms (Section~\ref{sec:ladder}).} We formalize the logarithmic peel
ladder as an algorithm with an exact two-level entry (Lemma~\ref{lem:entry}) and an
a-posteriori certificate (Proposition~\ref{prop:cert}): entering the ladder at a tower value
$\ge M$ incurs a total phase error $\le K d^{-M}$. Newton inversion of the phase map is
uniformly well conditioned since $\widetilde H'\in[0.9971,\,1.0030]$ across all measured
pairs. We also identify, quantify and turn into a guard rule the dominant instrument
systematic: truncation of the engine's Kneser theta series at $n_{\mathrm{lim}}$ harmonics
produces a deterministic mode floor at level $\approx e^{-2.95\,n_{\mathrm{lim}}}$ with onset
$k\approx 2.3\,n_{\mathrm{lim}}$, independent of working precision.

\emph{Regularity dictionary (Section~\ref{sec:dictionary}).} For branch pairs of finite
smoothness we prove the transport statement ``branch $C^m$ with seam jumps $\Rightarrow$
$\lvert\widehat P(k)\rvert\asymp k^{-(m+2)}$'' with an explicit two-seam interference
structure (Propositions~\ref{prop:fourierjump}, \ref{prop:seams}). On the explicit $C^1$
branches of \cite[Rem.~7.5(ii)]{unified} the measured exponents are $3.06$ and $3.006$,
calibrating the pipeline's ability to detect polynomial signatures down to amplitudes
$10^{-150}$.

\emph{Mode decay and the boundary singularity (Sections~\ref{sec:decay},
\ref{sec:singularity}).} Clean Fourier modes of the regular phases to $k=300$ show a decay
profile $\sigma_\eff(k)$ that is \emph{universal} across all fourteen measured pairs. A direct
complexified evaluation of the phase (a $y$-scan of the peel ladder) locates an isolated
singularity per period at $z_*\approx x_*+0.055\,i$; its strength and type are pair-universal
to $0.01$--$0.3\%$, only the abscissa $x_*(b,d)$ varies. Steepest-descent asymptotics
(Proposition~\ref{prop:saddle}) identify the measured four-parameter decay law
$\ln A_k = \mathrm{const}+\gamma\ln k-2\pi\sigma_* k-\sqrt{\pi c k}$ with an essential
singularity of exponential type on the strip boundary, and two elementary consequences---a
crossover scale (Lemma~\ref{lem:crossover}) and a digit-yield formula
(Corollary~\ref{cor:yield})---explain the observed fit-window drift and give the mode counts
needed for base-change tables.

\emph{First-order theory of amplitudes (Section~\ref{sec:kernel}).} Differentiating the family
of Abel functions in the base parameter $\beta=\ln\ln b$ yields a \emph{response kernel}
$V_d$ solving the linear cohomological equation
$V_d(d^{\,y})-V_d(y)=-A_d'(d^{\,y})\,y\,(\ln d)\,d^{\,y}$, which along the tower orbit
collapses to the strikingly simple increment $-T_d/T_d'$. Under a base-differentiability
hypothesis we prove $\Phi_{b,d}=\varepsilon\,G_d+O(\varepsilon^2)$ with
$\varepsilon=\ln\alpha$ and an explicitly $1$-periodic kernel profile $G_d$
(Theorem~\ref{thm:kernel}), which upgrades the empirical amplitude law
$A_k\approx C_k\lvert\ln\alpha\rvert$ to a derived first-order statement and identifies the
hub slope $\frac{d}{d(\ln\ln b)}\mu_{e,b}\big|_{b=e}=\int_0^1 G_e$, measured as
$2.191463277066$. Ultra-near-diagonal experiments determine $C_1(e)$ to seventeen digits.

\emph{The mean-shift hub (Section~\ref{sec:hub}).} The exact antisymmetry of \cite{unified}
makes the decomposition $\mu_{b,d}=m(d)-m(b)+\Delta_{b,c_*;d}$ exact with a single defect
integral; we derive the third-order defect formula and validate the chain
$\Delta^{(2)}+\Delta^{(3)}$ to $\sim10^{-15}$ absolute. The hub function $m(b)=\mu_{e,b}$ was
measured over $24$ bases in $[1.46,100]$; at the lower edge we \emph{derive} the divergence
law $m(b)\sim-\pi\sqrt{2\eta/e}\,(b-\eta)^{-1/2}$ from the parabolic bottleneck of the
threshold base (Lemma~\ref{lem:edge}), matching the fitted coefficient to $2.3\%$. PSLQ
searches for closed forms of $m(2),m(3),m(5),m(10)$ and $C_1(e)$ against an extended constant
basis were uniformly negative.

\emph{A hybrid data format (Section~\ref{sec:format}).} We specify the format in which these
results become an implementable base-change primitive: hub column, per-pair Fourier head plus
universal decay curve for the working range $\lesssim 150$ digits, certified peel ladder
beyond, and cross-base identities (mean antisymmetry to $10^{-38}$, groupoid residuals) as
proof residuals with thresholds derived from the certificate, not from experience.

\subsection{Standing conditionality and provenance}

All statements about the regular Kneser--Paulsen--Cowgill branch are conditional in exactly
the sense of \cite[Rem.~7.4]{unified}: they consume the real-channel and peeling-growth
hypotheses of \cite[Hyp.~2.2, 2.4]{phaselaw}, which remain to be derived from the branch
construction itself. The toy-branch results of Section~\ref{sec:dictionary} are
unconditional. Numerical values are quoted with the precision at which they passed the gates
of Appendix~\ref{app:repro}; the raw data, scripts and run metadata reside in the project
repository.

\section{Preliminaries}\label{sec:prelim}

We keep the notation of \cite{phaselaw,unified}. Throughout, $b,d>1$ are admissible bases,
$T_a$ denotes the regular increasing branch normalized by \eqref{eq:tetration} (frame
(B1)--(B2) of \cite{unified}), $A_a=\slog_a$ its Abel inverse, $L_b=\log_b$, and
\[
\alpha=\log_b d,\qquad
D_n(\theta)=D^{b,d}_n(\theta)=L_b^{\circ n}\bigl(T_d(n+\theta)\bigr),\qquad
\theta\in[0,1].
\]
Under \cite[Hyp.~2.2, 2.4]{phaselaw}, $D_n\to D$ uniformly, the phase law
\eqref{eq:phaselaw} holds with $\Phi_{b,d}=A_b\circ D-\mathrm{id}$, and
$s_j:=\inf_{\theta}T_d(j+\theta)$ satisfies the exact tower recursion $s_{j+1}=d^{\,s_j}$
(\cite[Lem.~2.5]{unified}), so $\sum_j 1/s_j<\infty$ automatically. The $\sigma$-coordinates
of \cite[Lem.~6.1, 6.3]{unified} write the finite tails as
\begin{equation}\label{eq:sigmarec}
D_n'(\theta)=\alpha\,T_d'(\theta)\prod_{j=1}^{n-1}\bigl(1+\sigma^{(n)}_j(\theta)\bigr),
\qquad
\sigma^{(n)}_j=\frac{L_b(\alpha)+L_b\bigl(1+\sigma^{(n)}_{j+1}\bigr)}{\alpha\,T_d(j+\theta)},
\qquad \sigma^{(n)}_{n-1}=0,
\end{equation}
with the uniform tail bound $\lvert\sigma^{(n)}_j\rvert\le C_0/s_j\le\tfrac12$ for
$j\ge J_0$, where $C_0=\bigl(\lvert L_b(\alpha)\rvert+\ln 2/\ln b\bigr)/\alpha$
(\cite[Eq.~(17)]{unified}). We write $\widehat P(k)$ for the complex Fourier coefficients of
the centered profile $P=P_{b,d}$, $A_k=2\lvert\widehat P(k)\rvert$ for the real mode
amplitudes, and define the \emph{local decay rate}
\begin{equation}\label{eq:sigmaeff}
\sigma_\eff(k)\;=\;\frac{\ln A_k-\ln A_{k+1}}{2\pi}.
\end{equation}
A strip-analytic $P$ with half-width $\sigma$ and no boundary structure would give
$\sigma_\eff(k)\to\sigma$; a $C^m$ profile with seam jumps gives
$\sigma_\eff(k)\approx\frac{m+2}{2\pi k}$.

For an actual finite height $x=n+\theta$ we write
$R_n(\theta):=F_{b,d}(n+\theta)-(n+\theta)-\Phi_{b,d}(\theta)$ for the
\emph{finite-height remainder}; $\lVert R_n\rVert_\infty\to0$ at inverse-tower speed but
$R_n\not\equiv0$, so any claim that a phase table evaluates $F_{b,d}$ at finite $x$ must
carry this term in its error budget. The stopped-coordinate certificate of Paper~III
\cite{paperIII} bounds it by $2C_0'/T_b(n-2)$, which jumps from $\sim5\times10^{-7}$ at
$n=5$ to $\sim10^{-1{,}650{,}000}$ at $n=6$ for $(e,2)$: the switch height below which the
direct engine must be used is tiny and pair-computable.

\section{Certified evaluation: the peel ladder}\label{sec:ladder}

The evaluation primitive behind every measurement in this paper is the logarithmic peel
ladder. Its correctness rests on two exact identities and one hyper-fast tail bound; we record
them in a form that yields a certificate per evaluation.

\begin{lemma}[Exact two-level entry]\label{lem:entry}
For every $m\ge 0$ and $\theta\in[0,1]$ for which the first two peels stay in the real
channel,
\begin{equation}\label{eq:entry}
L_b^{\circ 2}\bigl(T_d(m+2+\theta)\bigr)\;=\;\log_b\alpha+\alpha\,T_d(m+\theta).
\end{equation}
\end{lemma}

\begin{proof}
$L_b\bigl(T_d(m+2+\theta)\bigr)=L_b\bigl(d^{\,T_d(m+1+\theta)}\bigr)=\alpha\,T_d(m+1+\theta)$,
and then
$L_b\bigl(\alpha T_d(m+1+\theta)\bigr)=\log_b\alpha+L_b\bigl(d^{\,T_d(m+\theta)}\bigr)
=\log_b\alpha+\alpha\,T_d(m+\theta)$.
\end{proof}

Thus the ladder never needs a tower value beyond $T_d(m+\theta)$ in memory: the two topmost
levels are absorbed exactly, and everything above them is a hyper-exponentially small
$\sigma$-correction.

\begin{algorithm}[Certified peel ladder]\label{alg:ladder}
Given $(b,d)$, $\theta\in[0,1)$, a working threshold $M$ and a landing window
$\mathcal W=[w_-,w_+]\subset(0,\infty)$:
\begin{enumerate}
\item lift the tower from Taylor data on $[0,1]$ until $T_d(m+\theta)\ge M$;
\item enter with $W_0:=\log_b\alpha+\alpha\,T_d(m+\theta)$ (two levels absorbed exactly by
Lemma~\ref{lem:entry});
\item peel $W\leftarrow L_b(W)$ until $W\in\mathcal W$, counting $p$ peels;
\item return $\Phi(\theta)=A_b(W)+p-m-\theta$: since
$\slog_b\bigl(T_d(m+2+\theta)\bigr)=A_b(W)+(2+p)$ after the two absorbed levels and $p$
peels, subtracting the height $m+2+\theta$ collapses the bookkeeping, and a single
$\slog_b$ call on the landing window closes the evaluation.
\end{enumerate}
\end{algorithm}

\begin{proposition}[A-posteriori certificate]\label{prop:cert}
Assume the phase-law package and the quantitative channel of \cite{unified} for $(b,d)$
(by Paper~III \cite{paperIII}, the growth hypothesis and finite-chain positivity hold
unconditionally in stopped coordinates, and the limit channel reduces to one offset
inequality per pair), and
let $J_0$, $C_0$ be as in \eqref{eq:sigmarec}. If the ladder enters at a level $m$ with
$s:=T_d(m+\theta)\ge M\ge s_{J_0}$, then its output $\Phi^{(m)}(\theta)$ satisfies
\[
\bigl|\Phi^{(m)}(\theta)-\Phi_{b,d}(\theta)\bigr|\;\le\;K\,d^{-M},
\]
where $K$ depends only on $C_0$, on the Lipschitz constant of $A_b$ on the landing window,
and on the (finitely many) channel constants of the sub-window levels.
\end{proposition}

\begin{proof}[Proof sketch]
Entering at level $m$ is the same as evaluating the exact infinite object with all
corrections $\sigma_j$ at levels $j$ above the entry set to zero. By
\cite[Lem.~6.3]{unified} the neglected corrections satisfy
$\lvert\sigma_j\rvert\le C_0/s_j$, and by the tower recursion $s_{j+1}=d^{\,s_j}$ the first
neglected one already obeys $C_0/s_{m+1}=C_0\,d^{-s}\le C_0\,d^{-M}$, with the remaining ones
smaller than any fixed multiple of it. Transporting this perturbation down the ladder
multiplies it by the product of the peel derivatives $1/(W_k\ln b)$, which is bounded by a
constant on the compact channel below the entry, and the final $\slog_b$ is Lipschitz on the
landing window. Summing the geometric-in-tower tail gives the stated bound.
\end{proof}

In the campaign the certificate was accompanied by an operational gate: adding one further
explicit tower level changed no returned phase value by more than $10^{-120}$ at $M=10^4$,
consistent with the bound $C_0\,d^{-10^4}$.

\begin{proposition}[Conditioned inversion]\label{prop:newton}
Let $\delta:=\max_\theta\lvert\Phi'(\theta)\rvert\le 2\pi\sum_k k\,A_k$. Then
$\widetilde H=\mathrm{id}+\Phi$ is a $C^1$ diffeomorphism with
$\widetilde H'\in[1-\delta,1+\delta]$; the fixed-point iteration
$\theta_{n+1}=u-\Phi(\theta_n)$ is a \emph{global} contraction with factor
$\delta$ (Newton acceleration is optional after a standard basin check); and
the inverse phase requires no separate
computation: $\widetilde H^{-1}=\widetilde H_{d,b}$ exactly \cite[Thm.~A]{unified}.
The bounds $1-\delta$ are themselves certifiable from the mode table via
$\delta\le2\pi\sum_k kA_k$ plus the tail envelope.
\end{proposition}

Across all fourteen measured pairs, $\delta\le 3.0\times10^{-3}$, i.e.\
$\widetilde H'\in[0.9971,\,1.0030]$ (best pair $(e,3)$ with $\pm 2.3\times10^{-4}$, worst
$(2,10)$ and $(10,2)$ with $\pm 3.0\times10^{-3}$): the phase maps are essentially perfectly
conditioned.

\subsection{An instrument law: the theta-truncation floor}\label{subsec:floor}

The engine evaluates $T_a$ and $A_a$ through Kneser's construction with the $1$-periodic
theta mapping truncated at $n_{\mathrm{lim}}$ harmonics. This truncation is invisible at the
level of individual function values at ordinary precision, but it is \emph{not} invisible to
high Fourier modes of $\Phi$: it injects a deterministic, precision-independent periodic
perturbation of the branch, which propagates linearly into $\Phi$ through the derivative
transfer of \eqref{eq:product}. The campaign isolated this systematic cleanly:

\begin{observation}[Mode floor]\label{obs:floor}
With $n_{\mathrm{lim}}=30$ the mode amplitudes of $(e,2)$ freeze at
$A\approx 2.6\times10^{-44}$ with onset $k\approx 56$; with $n_{\mathrm{lim}}=60$ the floor
moves to $A\approx 9.9\times10^{-83}$ with onset $k\approx 137$. The floor level is
bit-identical between working precisions $\mathrm{dps}=140$ and $200$ (agreement
$10^{-25}$), hence a systematic of the branch series, not rounding and not the ladder. The
two measurements give the instrument law
\[
\text{floor}(n_{\mathrm{lim}})\;\approx\;\text{const}\cdot e^{-2.95\,n_{\mathrm{lim}}}
=\text{const}\cdot e^{-2\pi\cdot 0.47\,n_{\mathrm{lim}}},
\qquad
k_{\text{floor}}\;\approx\;2.3\,n_{\mathrm{lim}},
\]
the rate $e^{-2\pi\cdot0.47}$ per harmonic being the decay rate of the engine's Kneser theta
series at its effective evaluation height.
\end{observation}

The practical guard rule is immediate and is enforced in the pipeline: \emph{modes beyond
$2.3\,n_{\mathrm{lim}}$ are artifacts by definition}; for a target band $k\le k_{\max}$ one
must initialize with $n_{\mathrm{lim}}\ge k_{\max}/2.3$ and choose the working precision from
the floor law (both the rate and the measured prefactor, which sits about five orders below
the pure rate extrapolation). Initialization costs scale accordingly
($n_{\mathrm{lim}}=60$ at $\mathrm{dps}=200$: $56$\,s and $95$\,s for bases $2$ and $e$;
$n_{\mathrm{lim}}=120$: $239$\,s and $196$\,s; one-time, disk-cached). A third
calibration point at $n_{\mathrm{lim}}=240$ (floor onset $k\approx530$ against the
predicted $2.3\cdot240=552$) confirms the law across an eightfold range. We emphasize the
methodological point: this is the same failure mode---silent theta truncation---that
historically corrupted reference values in the underlying engine; here it was detected
\emph{before} contaminating any conclusion, quantified, and converted into an instrument
constant with a formula.

\section{The regularity dictionary and its calibration}\label{sec:dictionary}

Question (Q1) is only meaningful if the measurement pipeline can distinguish regularity
classes. The companion note provides the exact tool: by \cite[Rem.~7.2]{unified},
$D=T_b\circ\widetilde H$, so given branch regularity, regularity of $\Phi$ and of the
limiting tail are equivalent, and \emph{branch} regularity is transported into \emph{phase}
regularity with known bookkeeping. For finite smoothness this yields a quantitative
dictionary.

\begin{proposition}[Fourier asymptotics with seam jumps]\label{prop:fourierjump}
Let $f:\Sone\to\R$ be $C^m$ and piecewise $C^{m+2}$, with $f^{(m+1)}$ having jump
discontinuities $J_1,\dots,J_r$ at $\tau_1,\dots,\tau_r\in\Sone$. Then, as
$k\to\infty$,
\[
\widehat f(k)\;=\;\frac{1}{(2\pi i k)^{m+2}}\sum_{i=1}^{r}J_i\,e^{-2\pi i k\tau_i}
\;+\;o\!\left(k^{-(m+2)}\right).
\]
\end{proposition}

\begin{proof}
Integrate by parts $m+1$ times; the boundary terms vanish by continuity of
$f,\dots,f^{(m)}$ on the circle. One further integration by parts on each smooth piece
produces the jump sum plus $(2\pi i k)^{-(m+2)}\widehat{f^{(m+2)}}(k)$ interpreted piecewise;
the latter is $o(1)$ by Riemann--Lebesgue.
\end{proof}

\begin{proposition}[Seam locations of the phase]\label{prop:seams}
Let the branches $T_b$, $T_d$ be piecewise real-analytic with $C^m$ seams at integer
arguments (as for the explicit branches of \cite[Rem.~7.5(ii)]{unified}). Then
$\Phi_{b,d}\in C^m(\Sone)$ is piecewise real-analytic, and within one period its seam set is
\[
\Sigma\;=\;\{0\}\;\cup\;\bigl\{\theta\in[0,1):\ \widetilde H_{b,d}(\theta)\in\Z\bigr\},
\]
i.e.\ the inner-branch seam at $\theta=0$ together with the finitely many points where the
phase map crosses an integer height of the outer branch. Consequently
$\widehat P(k)=O(k^{-(m+2)})$, with the leading envelope governed by the finite
seam sum of Proposition~\ref{prop:fourierjump} (two-sided bounds $\asymp$ require a
non-cancellation condition on that trigonometric sum) and with an interference oscillation between the
seam contributions of Proposition~\ref{prop:fourierjump}, of beat period
$\Delta k\approx 1/\mathrm{dist}(\Sigma)$ in $k$.
\end{proposition}

\begin{proof}[Proof sketch]
For each $\theta$, Algorithm~\ref{alg:ladder} expresses $\Phi(\theta)$ through finitely many
maps: the Taylor lift of $T_d$ on $[0,1]$ (seams only where $m+\theta$ crosses $\Z$, i.e.\
$\theta=0$), real logarithms (analytic on $(0,\infty)$), and one evaluation of $A_b$ (seams
where its argument crosses the seam values $T_b(\Z)$, i.e.\ where
$A_b(D(\theta))=\widetilde H(\theta)\in\Z$). Away from $\Sigma$ all maps are analytic; the
$C^m$ matching at the seams is the branch matching, transported by the chain rule; the seam
count is finite because $\widetilde H$ is a homeomorphism of degree one. The Fourier statement
is Proposition~\ref{prop:fourierjump}.
\end{proof}

\paragraph{Calibration measurement.}
On the explicit $C^1$ branches ($m=1$; quadratic critical piece, seams with $T''$ jumps at
integers) the dictionary predicts $\lvert\widehat P(k)\rvert\asymp k^{-3}$. Log--log fits over
$k\in[4,60]$ give decay exponents $p=3.06$ for both orientations of $(e,2)$ and $p=3.006$ for
$(2,10)$, with the predicted interference oscillation clearly visible; for the toy pair
$(e,2)$, $\mu_{\mathrm{toy}}\approx-1.1290$ places the outer seam at
$\theta\approx0.1290$, so the predicted beat period is $\Delta k\approx 1/0.129\approx7.8$,
consistent with the observed oscillation. Two further pipeline gates were passed on the toy
branches: $\mu_{\mathrm{toy}}(e,2)=-1.1290411574965$ agrees with the independent computation
of \cite[Rem.~7.5]{unified} up to $-6.6\times10^{-11}$, which is exactly the expected
$N^{-3}$ aliasing error of the $512$-point grid; and the inverse identity
$\Phi_{2,e}(\widetilde H_{e,2}(\theta))+\Phi_{e,2}(\theta)=0$ holds to $5.7\times10^{-101}$.
The pipeline therefore provably detects polynomial signatures; every regularity claim below
is calibrated against this control.

\section{Mode decay of the regular phases}\label{sec:decay}

For the regular branch the picture is entirely different from the toy control: no polynomial
signature appears anywhere down to amplitudes $10^{-150}$, and no geometric regime appears
either.

\begin{observation}[Concave decay, universal profile]\label{obs:decay}
For $(e,2)$ at $n_{\mathrm{lim}}=120$, $\mathrm{dps}=190$, $N=1024$ grid points, the modes are
clean to $k\approx290$ (no floor visible), with local rates
\[
\begin{array}{c|ccccccccccc}
k & 1 & 2 & 3 & 5 & 8 & 20 & 40 & 60 & 100 & 200 & 290\\\hline
\sigma_\eff(k) & 0.740 & 0.567 & 0.487 & 0.404 & 0.344 & 0.260 & 0.2154 & 0.1946 & 0.1726 &
0.1484 & 0.1375
\end{array}
\]
monotone and strictly concave in $\ln k$ throughout. The profiles $\sigma_\eff(k)$ for
$k\gtrsim5$ agree to $3$--$4$ decimals across \emph{all} measured pairs---$(e,2)$, $(2,e)$,
$(2,10)$, $(10,2)$, $(e,10)$, $(10,e)$, $(e,3)$, $(3,e)$, $(e,5)$, $(5,e)$, $(3,5)$, $(5,3)$,
$(1.5,e)$, $(e,1.5)$: the decay \emph{shape} is a universal function $\Xi(k)$, and the pair
enters essentially only through the amplitude $C_1\lvert\ln\alpha\rvert$ (plus a mild shape
drift $e^{\delta k}$ with $\lvert\delta\rvert\sim10^{-3}$).
\end{observation}

Purely real-axis data to $k=290$ admit two competing tail models: a strip of positive
half-width $\sigma_\infty\in[0.05,0.09]$ approached as $\sigma_\eff=\sigma_\infty+\beta/\sqrt
k$, or a stretched-exponential $\ln A_k\sim-ck^{0.79}$ with $\sigma_\infty=0$, i.e.\ a Gevrey
class far beyond all finite-smoothness classes. The first wins on $k\in[20,140]$, the second
on $k\in[140,290]$; no two-parameter model fits globally. The dichotomy is resolved not by
more modes but by going complex, which is the subject of the next section. What is already
decided at this stage: \cite[Conj.~8.1]{phaselaw} \emph{cannot} hold in its naive form with
the first-mode rate ($\sigma\approx0.74$)---the strip, if any, is an order of magnitude
narrower---and the phase is not of any finite smoothness class.

For applications the dispute is immaterial in the working range: the measured digit yield
$D(k)=\log_{10}(1/A_k)$ is
\[
D(60)=47,\qquad D(100)=67,\qquad D(120)=76,\qquad D(300)\approx148,
\]
identical for both models where measured, and both extrapolate $500$ digits to
$k\approx1300$--$1700$ modes (Corollary~\ref{cor:yield}).

\section{The boundary singularity}\label{sec:singularity}

\subsection{Direct complexified evaluation}

The peel ladder complexifies: for $\theta+iy$ with $y>0$ the tower, the peels and the final
$\slog_b$ all extend, as long as every intermediate value avoids the logarithmic cut
$(-\infty,0]$ and the solver tracks a consistent branch. The campaign implemented this as a
bottom-up solver for $\mathrm{sexp}_b(J+z+\varphi)=\alpha\,T_d(J+z)\cdot(1+\sigma\text{-closure})$
with the entry level $J$ chosen per point, secant iteration with damping, continuation along
paths in $y$ and $x$ (never cold evaluation near the critical zone), and, crucially,
\emph{mid-strip validation}: away from the critical zone the solver agrees with the
$k\le300$ Fourier reconstruction to below the per-point closure bounds (typically
$10^{-20}$--$10^{-60}$), so both instruments certify each other where both are valid.

\begin{observation}[Breach geometry]\label{obs:breach}
For $(e,2)$ the solver is regular on the whole strip up to $y\approx0.16$ \emph{except} in a
single slanted sector per period: at $y=0.060$ the failure zone is
$x\in[1.004,1.022]$, widening through $[0.982,1.012]$ at $y=0.065$, $[0.964,1.000]$ at
$0.070$, $[0.91,0.98]$ at $0.085$, $[0.879,\ge0.95]$ at $0.095$, with center slope
$dx/dy\approx-2.8$ and apex
\[
z_*\;\approx\;(1.02\pm0.02)\;+\;(0.055\pm0.005)\,i \pmod 1 .
\]
Inside the zone the level-$J$ system exhibits quantized sheet slips and non-convergence
(root density), with closure bounds down to $10^{-305}$---excluding any truncation artifact.
Off the zone, lines continue smoothly far beyond the apex height (measured cleanly to
$y\approx0.164$), with indications of a second structure at $y\approx0.17$--$0.22$,
$x\approx0.6$--$0.7$, left to future work. The mode phases corroborate the location: the
drift of $\arg\widehat P(k)$ defines $x_*(k)$ moving $0.47\to0.65$ over $k\le280$ with a
$1/\sqrt k$ approach toward $x_*\approx0.94$--$1.0$.
\end{observation}

\begin{observation}[Universality of the singularity]\label{obs:universal}
At fixed height $y=0.085$ the failure zones of three structurally different pairs are
congruent: $(e,2)$: $[0.91,0.98]$; $(3,5)$: $[1.01,1.08]$ (a rigid shift by $+0.100$, exactly
the phase-fit shift); $(2,10)$: $[0.650,0.710]$---same height, same width
($\approx0.06$--$0.07$), same shape; only the abscissa $x_*(b,d)$ is pair-dependent
(centers $\approx0.945/1.045/0.680$; the near-round values invite scrutiny). Structure fits
over $220$ modes (amplitudes and phases, $k\in[60,280]$) with the one-singularity saddle
model of Proposition~\ref{prop:saddle} give
\begin{align*}
(e,2):&\ \ \sigma_*=0.0789,\quad c=59.78,\quad \gamma=9.47\quad(\mathrm{rms}\ 0.005\ \text{over
220 modes});\\
(3,5):&\ \ \sigma_*=0.0791,\quad c=59.77,\quad \gamma=9.59,
\end{align*}
i.e.\ $\sigma_*$ pair-equal to $2.4\times10^{-4}$ and $c$ to $0.013\%$; the phase fit yields
in addition a complex strength with $\arg$-parameter $\chi\approx0.68$ rad, pair-universal to
$0.3\%$.
\end{observation}

\subsection{Saddle-point asymptotics}

The measured four-parameter law is exactly the Fourier signature of an isolated essential
singularity of exponential type sitting on the strip boundary.

\begin{proposition}[Singularity $\Rightarrow$ decay law]\label{prop:saddle}
Let $P$ be real on $\R$, $1$-periodic, holomorphic on
$\{\lvert\operatorname{Im}\theta\rvert<\sigma_*\}$ except for isolated singularities per
period at $\theta_*=x_*+i\sigma_*$ and $\bar\theta_*$, where locally
\[
P(\theta)\;=\;h(\theta)\,\exp\!\Bigl(\frac{v}{\theta-\bar\theta_*}\Bigr),\qquad
h(\theta)\sim h_0\,(\theta-\bar\theta_*)^{-\rho},
\]
for some $v\in\mathbb C\setminus\{0\}$, $\rho\in\R$, and suppose the contour can be deformed
to the boundary away from the singularity. Then, as $k\to+\infty$,
\[
\ln\bigl|\widehat P(k)\bigr|\;=\;\mathrm{const}+\gamma\ln k-2\pi\sigma_*k-\sqrt{\pi c\,k}
+o(\sqrt k),
\qquad
\arg\widehat P(k)\;=\;-2\pi k\,x_*+O(\sqrt k),
\]
where the strength satisfies $8\lvert v\rvert\cos^2\Psi=c$ with
$\Psi$ the steepest-descent phase determined by $\arg v$, and
$\gamma=\rho/2-3/4$ collects the algebraic prefactor ($h\sim h_0 s^{-\rho}$) and the
$O(k^{-3/4})$ saddle-width contribution. In particular, for real negative $v=-a$ one has
$\Psi=\pi/4$ and $a=c/4$.
\end{proposition}

\begin{proof}[Proof sketch]
For $k>0$ push the contour of $\widehat P(k)=\int_0^1P(\theta)e^{-2\pi ik\theta}\,d\theta$ to
$\operatorname{Im}\theta=-\sigma_*$; the smooth part contributes
$O(e^{-2\pi\sigma_*k})$ with bounded prefactor, and the dominant term comes from the
neighborhood of $\bar\theta_*$. Writing $\theta=\bar\theta_*+s$, the local integral is
$e^{-2\pi ik\bar\theta_*}\int h\,\exp\bigl(v/s-2\pi iks\bigr)\,ds$ with saddle
$s_*=\sqrt{iv/(2\pi k)}$ and saddle value $\varphi(s_*)=-4\pi iks_*$, so
$\lvert e^{\varphi(s_*)}\rvert=\exp\bigl(-\sqrt{8\pi\lvert v\rvert k}\,\cos\Psi\bigr)$ with
$\cos\Psi$ fixed by the admissible steepest-descent branch of $\sqrt{iv}$; the saddle width
$\lvert s_*\rvert\propto k^{-1/2}$ and the prefactor $h_0 s_*^{-\rho}$ produce the algebraic
factor $k^{\gamma}$. Matching $\sqrt{8\pi\lvert v\rvert}\cos\Psi=\sqrt{\pi c}$ gives the
strength relation; standard steepest-descent estimates (e.g.\ \cite[Ch.~4]{olver}) control
the error terms. The check for real $v=-a$: the saddle direction gives
$\cos\Psi=1/\sqrt2$ and $\lvert e^{\varphi(s_*)}\rvert=e^{-2\sqrt{\pi a k}}$, i.e.\ $a=c/4$.
\end{proof}

With the fitted $c\approx59.8$ this puts the singularity strength at
$\lvert v\rvert\cos^2\Psi=c/8\approx7.5$, with $\arg v$ determined by the measured
$\chi\approx0.68$ rad. Two elementary consequences organize the real-axis phenomenology
completely.

\begin{lemma}[Crossover scale]\label{lem:crossover}
In the model of Proposition~\ref{prop:saddle} the increments of $-\ln A_k$ are dominated by
the $\sqrt k$ term for $k<k_c$ and by the linear term for $k>k_c$, where
\[
k_c\;=\;\frac{c}{16\pi\sigma_*^2}.
\]
\end{lemma}

\begin{proof}
Equate the derivatives $2\pi\sigma_*$ and $\tfrac12\sqrt{\pi c/k}$.
\end{proof}

For $c=59.8$: $k_c\approx190$ if $\sigma_*=0.079$, but $k_c\approx390$ if
$\sigma_*=0.055$. This single number explains the entire fit phenomenology: window fits below
$k_c$ systematically \emph{overestimate} $\sigma_*$ (the $\sqrt k$ term masquerades as a
larger linear rate), which is exactly the observed downward drift of subwindow estimates
($0.0822\to0.0746$ as the window moves up), converging toward the directly measured apex
height $0.055\pm0.005$. The discrimination experiment was subsequently run: an
$n_{\mathrm{lim}}=240$ campaign ($\mathrm{dps}=260$, $N=2048$) yields clean modes to
$k\approx520$ (floor onset $\approx530$, a third confirmation of the instrument law of
Observation~\ref{obs:floor}), and the sequence of sliding three-parameter window fits
converges with a clean $1/w$ law, $\sigma_*(w)=0.062+2.83/w$ reproducing all windows to
$\pm5\times10^{-4}$, extrapolating to $\sigma_*^{\infty}=0.062\pm0.003$. Note that no
single mode value discriminates (holding the subleading coefficients fixed misreads
$\sigma_\eff(500)=0.1237$ as $\sigma_*\approx0.078$); only the window-sequence
extrapolation separates the hypotheses. Combining with the apex,
\[
\boxed{\ \sigma_*\;=\;0.058\pm0.006\ }
\]
and the fitted $0.079$ is \emph{excluded} as an asymptotic value: it is the
pre-asymptotic effective rate of windows below the crossover.

\begin{corollary}[Digit yield and mode counts]\label{cor:yield}
In the fitted model the digit yield $D(k)=\log_{10}(1/A_k)$ satisfies
$D(k)\ln10=-\ln C-\gamma\ln k+2\pi\sigma_*k+\sqrt{\pi ck}$. Calibrating the single constant
$\ln C$ at the measured $D(300)=148$ reproduces the measured yields to $\pm1$ digit
($D(60)=46.0$ vs.\ $47$; $D(100)=65.9$ vs.\ $67$; $D(120)=75.1$ vs.\ $76$) and gives
$D=500$ at $k\approx1400$ (central value; $1300$--$1700$ across the admissible parameter
range), i.e.\ $n_{\mathrm{lim}}\approx600$--$750$ with a one-time initialization of
$1$--$2$ hours per base.
\end{corollary}

\begin{conjecture}[Universal singularity data]\label{conj:universal}
For all admissible pairs $(b,d)$, $b,d>e^{1/e}$, the phase $\Phi_{b,d}$ extends
holomorphically to the punctured strip
$\{\lvert\operatorname{Im}\theta\rvert<\sigma_*\}\setminus(\{\theta_*\}+\Z)$ with an isolated
essential singularity of exponential type at $\theta_*=x_*(b,d)+i\sigma_*$; the height
$\sigma_*$ is independent of the pair; the normalized strength $\lvert v\rvert$ and phase
$\arg v$ are pair-independent to leading order (the diagonal kernels of
Section~\ref{sec:kernel} carry a smooth inner-base dressing of $0.4\%$ at $k=2$ growing to
$\sim9\%$ at $k=8$, so exact universality can hold at most for the height and the
singularity type),
and $x_*(b,d)$ equals the accumulation point of the mode phases,
$x_*=-\lim_k \arg\widehat P(k)/(2\pi k)$ along the drift-corrected sequence. In particular
\cite[Conj.~8.1]{phaselaw} holds in the punctured form, with universal boundary data.
\end{conjecture}

The universality of the profile $\Xi(k)$ (Observation~\ref{obs:decay}) and of the fitted
$(\sigma_*,c,\chi)$ (Observation~\ref{obs:universal}) is strong evidence; the linear-response
theory of the next section reduces the conjecture, for near-diagonal pairs, to a statement
about a single one-parameter family of kernels.

\section{First-order theory: the response kernel}\label{sec:kernel}

The empirical amplitude law $A_k(b,d)\approx C_k\,\lvert\ln\alpha\rvert$, with $C_1$ constant
to $\pm2.2\%$ over the full measured range $\ln\ln b\in[-0.90,0.83]$, calls for perturbation
theory in the mismatch parameter
\[
\varepsilon\;:=\;\ln\alpha\;=\;\ln\ln d-\ln\ln b,
\]
which vanishes exactly on the diagonal $b=d$, where $\Phi_{d,d}\equiv0$. The natural base
coordinate is $\beta:=\ln\ln b$.

\begin{hypothesis}[Base differentiability]\label{hyp:basediff}
The family $b\mapsto(T_b,A_b)$ of regular branches is $C^1$ in $\beta$ near $\beta_0=\ln\ln
d$, locally uniformly on compacts of the respective domains, with continuous derivative
$V_d(y):=\partial_\beta A_b(y)\big|_{b=d}$; moreover
$\partial_\beta\Phi_{b,d}(\theta)\big|_{b=d}=\lim_n V_d\bigl(T_d(n+\theta)\bigr)$ with the
interchange of limit and derivative uniform in $\theta$.
\end{hypothesis}

For the intended branches this is supported by Paulsen's holomorphy of the construction in
the base parameter \cite{paulsen} and is numerically evident; we state it as a hypothesis in
the same conditional spirit as \cite[Rem.~7.4]{unified}.

\begin{lemma}[Cohomological equation]\label{lem:cohom}
Under Hypothesis~\ref{hyp:basediff}, differentiating the global Abel equation
$A_b(b^{\,y})=A_b(y)+1$ in $\beta$ at $b=d$ gives, for all $y\in\R$,
\begin{equation}\label{eq:cohom}
V_d(d^{\,y})-V_d(y)\;=\;-A_d'(d^{\,y})\;y\,(\ln d)\,d^{\,y},
\end{equation}
and along the tower orbit $y=T_d(x)$ the right-hand side collapses:
\begin{equation}\label{eq:orbit}
V_d\bigl(T_d(x+1)\bigr)-V_d\bigl(T_d(x)\bigr)\;=\;-\,\frac{T_d(x)}{T_d'(x)} .
\end{equation}
\end{lemma}

\begin{proof}
$\partial_\beta b^{\,y}=y(\ln b)b^{\,y}$ since $d\ln b/d\beta=\ln b$; evaluate at $b=d$ to
get \eqref{eq:cohom}. For \eqref{eq:orbit} use $A_d'(z)=1/T_d'(A_d(z))$, $A_d(d^{\,y})=x+1$,
and $T_d'(x+1)=(\ln d)T_d(x+1)T_d'(x)$:
$A_d'(d^{\,y})\,y(\ln d)d^{\,y}
=\dfrac{T_d(x)\,(\ln d)\,T_d(x+1)}{(\ln d)\,T_d(x+1)\,T_d'(x)}=\dfrac{T_d(x)}{T_d'(x)}$.
\end{proof}

Since $T_d'(n)\gtrsim(\ln d)^n\prod_{i<n}T_d(i)$ by iterating the derivative transfer, the
increments \eqref{eq:orbit} decay at inverse-tower speed, and the following limit is
hyper-fast convergent.

\begin{theorem}[Response kernel and first-order phase]\label{thm:kernel}
Assume the phase-law package for pairs near the diagonal and
Hypothesis~\ref{hyp:basediff}. Define
\[
S_d(\theta)\;:=\;\sum_{j\ge0}\frac{T_d(j+\theta)}{T_d'(j+\theta)}
\qquad(\text{hyper-exponentially convergent}),
\qquad
G_d(\theta)\;:=\;S_d(\theta)-V_d\bigl(T_d(\theta)\bigr).
\]
Then $G_d$ is exactly $1$-periodic and continuous, and as $b\to d$,
\[
\Phi_{b,d}(\theta)\;=\;\varepsilon\,G_d(\theta)\;+\;O(\varepsilon^2)
\qquad\text{uniformly in }\theta .
\]
\end{theorem}

\begin{proof}
Write $v(x):=V_d(T_d(x))$. By \eqref{eq:orbit},
$v(n+\theta)=v(\theta)-\sum_{j=0}^{n-1}T_d(j+\theta)/T_d'(j+\theta)$, so
$\lim_n v(n+\theta)=v(\theta)-S_d(\theta)=-G_d(\theta)$ exists with hyper-fast convergence.
By Hypothesis~\ref{hyp:basediff},
$\partial_\beta\Phi_{b,d}(\theta)\big|_{b=d}=\lim_n v(n+\theta)=-G_d(\theta)$, and since
$\varepsilon=\mathrm{const}-\beta$ at fixed $d$, $\partial_\varepsilon\Phi\big|_{b=d}=G_d$;
Taylor's theorem with the locally uniform derivative gives the expansion. Periodicity:
$G_d(\theta+1)-G_d(\theta)=\bigl(S_d(\theta+1)-S_d(\theta)\bigr)
-\bigl(v(\theta+1)-v(\theta)\bigr)
=-\tfrac{T_d(\theta)}{T_d'(\theta)}-\bigl(-\tfrac{T_d(\theta)}{T_d'(\theta)}\bigr)=0$,
using \eqref{eq:orbit} once more.
\end{proof}

\begin{corollary}[First-order laws]\label{cor:firstorder}
With $\widehat G_d(k)$ the Fourier coefficients of $G_d$:
\begin{itemize}
\item[(i)] \emph{Amplitude law.} $A_k(b,d)=2\lvert\widehat
G_d(k)\rvert\,\lvert\varepsilon\rvert+O(\varepsilon^2)$; hence
$C_k(d)=2\lvert\widehat G_d(k)\rvert$, and the observed near-universality of $C_1$ in $d$ is
the statement that $d\mapsto2\lvert\widehat G_d(1)\rvert$ is nearly flat.
\item[(ii)] \emph{Mean law and hub slope.}
$\mu_{b,d}=\varepsilon\,\overline{G}_d+O(\varepsilon^2)$ with
$\overline G_d=\int_0^1G_d$; combined with the exact antisymmetry \cite[Thm.~B]{unified},
the hub function $m(b)=\mu_{e,b}$ satisfies
\[
\frac{dm}{d(\ln\ln b)}\Big|_{b=e}\;=\;\overline G_e .
\]
\item[(iii)] \emph{Universal diagonal shape.} The centered profile shape at the diagonal is
$G_d-\overline G_d$; its mode ratios are the $\varepsilon\to0$ limits of the measured pair
ratios.
\end{itemize}
\end{corollary}

\paragraph{Measurements.}
Ultra-near-diagonal runs $(d',e)$ with $\ln d'=1+\varepsilon$,
$\varepsilon\in\{10^{-12},10^{-15}\}$, extract the first-mode ripple ($\sim4\times10^{-19}$)
cleanly above the engine floor and determine
\begin{align*}
C_1(e)&\;=\;2\lvert\widehat G_e(1)\rvert\;=\;0.00038859729244516473\quad(\pm\text{1 ulp at
17 digits}),\\
\overline G_e&\;=\;2.191463277066,
\end{align*}
the latter measured as the hub slope $dm/d(\ln\ln b)\vert_{b=e}$; the diagonal shape ratios
are $A_2/A_1=0.0093355030$, $A_3/A_1=2.66505\times10^{-4}$, $A_4/A_1=1.26473\times10^{-5}$,
and the diagonal first-mode phase is $\varphi_1=1.33536363153$ rad---close to, but clearly
distinct from, $\operatorname{Im}L_e=1.33724$ ($L_e$ the principal fixed point of
$e^z$). Near-diagonal runs at $\pm0.02$ give the kernel amplitude curve
\[
C_1(d)\ [\times10^{-4}]:\quad 3.953\ (d{=}2),\quad 3.886\ (e),\quad 3.868\ (3),\quad
3.793\ (5),\quad 3.720\ (10),
\]
nearly affine in $\ln\ln d$ (slope $\approx-0.19\times10^{-4}$), with the quadratic
($\varepsilon^2$) deviations resolved at $\sim9\times10^{-6}\,\varepsilon$. These values were subsequently \emph{derived}: the selection principle for
$V_d$ is formulated, proved unique, and proved solvable in Paper~V of this series
\cite{paperV}, and its constructive Riemann--Hilbert realization reproduces the kernel
family engine-free to $30$--$34$ digits for $d\in\{2,e,3,5,10\}$, in mutual agreement with
$\varepsilon$-stencil differentiation of the branch family. Two independent routes thereby
exposed an error in a third (complex-step references at the extreme bases $d=2$ and $d=10$,
reliable only to $\approx20$ resp.\ $\approx15$ digits: the engine floor at
$\lvert\lambda(2)\rvert=1.23$ near the parabolic edge, and the narrowed theta strip at
$\ln 10$); all kernel values quoted here are confirmed by the surviving routes. The
second-order (in $\varepsilon$) kernels are measured as well, e.g.\
$\overline H_e=1.66735420625407$ with the $\ln\ln$-exact node parametrization, and the
two-term Taylor law $\widehat P_k=\varepsilon\widehat G_k+\varepsilon^2\widehat H_k$
reproduces the mode heads of \emph{all} measured pairs to third-order-small residuals even
at $\varepsilon=1.20$.

\section{The mean-shift hub}\label{sec:hub}

\subsection{Exact decomposition and higher-order defects}

\begin{proposition}[Exact hub]\label{prop:hub}
Fix a hub base $c_*$ and set $m(b):=\mu_{c_*,b}$. Then for all admissible $b,d$,
\[
\mu_{b,d}\;=\;m(d)-m(b)\;+\;\Delta_{b,c_*;d},
\qquad
\Delta_{b,c_*;d}\;=\;\int_0^1P_{b,c_*}\bigl(\{\theta+\Phi_{c_*,d}(\theta)\}\bigr)\,d\theta,
\]
exactly: the identity combines the integrated cocycle \cite[Eq.~(39)--(40)]{phaselaw} with
the exact antisymmetry $\mu_{b,c_*}=-\mu_{c_*,b}$ \cite[Thm.~B]{unified}. Expanding the
defect in the small profiles,
\[
\Delta^{(2)}=\sum_{k\ne0}2\pi ik\,\widehat P_{b,c_*}(k)\,e^{2\pi ik\mu_{c_*,d}}\,
\widehat P_{c_*,d}(-k),
\quad
\Delta^{(3)}=\frac12\sum_{k\ne0}(2\pi ik)^2\,\widehat P_{b,c_*}(k)\,
e^{2\pi ik\mu_{c_*,d}}\,\widehat{P_{c_*,d}^{2}}(-k),
\]
with $\bigl|\Delta-\Delta^{(2)}-\Delta^{(3)}\bigr|\le\tfrac16\,
\lVert P_{b,c_*}'''\rVert_\infty\,\lVert P_{c_*,d}\rVert_\infty^3$, i.e.\ the error carries
two more powers of the (tiny) amplitudes than $\Delta^{(2)}$ itself.
\end{proposition}

\begin{proof}
Exactness is the cited combination. For the expansion, Taylor
$P_{b,c_*}(\theta+\mu+P_{c_*,d})=P_{b,c_*}(\theta+\mu)+P_{b,c_*}'(\theta+\mu)P_{c_*,d}
+\tfrac12P_{b,c_*}''(\theta+\mu)P_{c_*,d}^2+O(\lVert P\rVert^3\lVert
P'''\rVert)$, integrate, note the first term has mean zero, and convert the two integrals by
Parseval, $\int fg=\sum\widehat f(k)\widehat g(-k)$, with
$\widehat{f'}(k)=2\pi ik\widehat f(k)$ and the shift factor $e^{2\pi ik\mu}$.
\end{proof}

The hub economics are the point: $N$ bases require one hub column of $N$ values (plus
low-$k$ mode heads) instead of $N^2$ pair runs. With $c_*=e$ the campaign validated the chain
at the level the formulas predict: the ratios $\Delta^{(2)}/\Delta_{\text{measured}}$ are
$0.999816$, $0.999900$, $0.999651$, $1.000074$ for $(2,10)$, $(10,2)$, $(3,5)$, $(5,3)$
(residual $=$ the $\Delta^{(3)}$ term), and including $\Delta^{(3)}$ the ratios become
$1.0027$, $1.0010$, $1.0005$, $0.9999$ respectively---the truncated hub chain
$\mu=m(d)-m(b)+\Delta^{(2)}+\Delta^{(3)}$ has the scale the remainder bound predicts: from
the displayed ratios, the largest tested defect $(2,10)$ closes to about
$1.7\times10^{-11}$ absolute and the closer pair $(3,5)$ to about $1.5\times10^{-13}$;
the full-precision chain in the repository carries the reconstruction further. As a
scaling check, $\Delta(3,5)=-2.96\times10^{-10}$ is $20\times$ smaller than
$\Delta(2,10)=6.4\times10^{-9}$, as befits the closer base pair.

\subsection{The hub function: values, edge law, large bases}

High-precision hub values (regular branch; $150$--$180$ digits available for $m(2)$, $\sim170$
for $m(3),m(5),m(10)$; leading digits shown):
\[
\begin{aligned}
m(1.5)&=-9.776277436687167726082341422480654566832\ldots\\
m(2)&=-1.128403776628924009879217902036426267112\ldots\\
m(3)&=+0.1924222937224320832082444066892813277903\ldots\\
m(5)&=+0.7712572961544758763804410539746626381043\ldots\\
m(10)&=+1.136557055713326440899375967586836969056\ldots
\end{aligned}
\]
together with a $24$-point raster over $b\in[1.46,100]$ at $40$ digits. (The value of
$\mu_{e,2}$ corrects the twelve-digit Table~1 entry of \cite{phaselaw}, whose known
$9.3\times10^{-9}$ error is exactly the inverse-sum residual reported there; the regenerated
table appears in Appendix~\ref{app:tables}. Pairwise mean antisymmetry now holds to
$\le10^{-39}$.)

At the lower edge the hub diverges, and the divergence law can be derived: it is the passage
time through the parabolic bottleneck of the threshold base.

\begin{lemma}[Parabolic edge law]\label{lem:edge}
Let $\eta=e^{1/e}$ and $b\downarrow\eta$. To leading order in the continuum (Fatou
coordinate) approximation of the bottleneck,
\[
m(b)\;=\;\mu_{e,b}\;=\;-\,\pi\sqrt{\frac{2\eta}{e}}\;(b-\eta)^{-1/2}\;
\bigl(1+o(1)\bigr),
\qquad
\pi\sqrt{2\eta/e}\;=\;3.23907\ldots,
\]
with a subleading correction of order $\ln(b-\eta)$.
\end{lemma}

\begin{proof}[Derivation]
At $b=\eta$ the map $t\mapsto b^{\,t}$ has the parabolic fixed point $t=e$ with multiplier
$1$. Write $b=\eta e^{\delta}$, $\delta=\ln(b/\eta)=(b-\eta)/\eta\,(1+O(b-\eta))$, and
$t=e+w$. Then
\[
b^{\,t}=\exp\bigl((e+w)(1/e+\delta)\bigr)
=e\cdot\exp\bigl(w/e+e\delta+w\delta\bigr)
=e+w+e^2\delta+\frac{w^2}{2e}+O(w\delta,w^3),
\]
so the orbit obeys $w_{n+1}=w_n+a+Bw_n^2+\cdots$ with $a=e^2\delta$, $B=1/(2e)$. The
continuum crossing time of such a bottleneck is
$\int_{-\infty}^{\infty}dw/(a+Bw^2)=\pi/\sqrt{aB}=\pi\sqrt{2/(e\delta)}
=\pi\sqrt{2\eta/e}\,(b-\eta)^{-1/2}(1+o(1))$. A $b$-tower therefore lags a comparison
$e$-tower by this many height units while traversing values near $e$, and the lag is exactly
what $\mu_{e,b}=\lim(\slog_e T_b-\mathrm{id})$ measures (with the sign of a delay). The
logarithmic subleading term is the standard second-order correction of parabolic passage
(Fatou coordinates; see \cite[\S10]{milnor}), and the derivation is made rigorous by the same
machinery; we use only the leading order here.
\end{proof}

The measured fit on $b\in[1.46,1.6]$ is
$m(b)=-A(b-\eta)^{-1/2}+B\ln(b-\eta)+C$ with $A=3.3148$, $B=-0.398$, $C=3.163$ (rms
$0.004$): the ratio $A/\pi\sqrt{2\eta/e}=1.023$ confirms the derived coefficient at the
$2.3\%$ level, the discrepancy being consistent with the fit window not yet asymptotic (the
prediction is that $A\to3.2391$ as the window closes on $\eta$; bases below $1.46$ were not
yet initialized). Notably the construction remains numerically robust arbitrarily close to
the edge: base $1.5$ (only $3.8\%$ above $\eta$) initializes in $1.2$ s and shows the same
universal decay profile, with $C_1=3.97\times10^{-4}$ drifting mildly above the bulk band.
At the other end, $m(b)-\ln\ln b$ decreases from $0.31$ at $b=8$ to $0.073$ at $b=100$; the
natural conjecture $m(b)=\ln\ln b+o(1)$ (with rate $O(1/\ln b)$?) is left open.

\subsection{Closed forms: negative evidence}

PSLQ searches at $150$ digits and coefficient bounds $10^6$--$10^8$ found \emph{no} integer
relations for $m(2)$ against $\{1,\ln2,1/\ln2,\ln\ln2,\ln2/(1-\ln2),e,\gamma\}$, nor for
$m(2),m(3),m(5),m(10)$ against an extended basis including $\eta$, $\pi$, $e$, logarithms,
iterated logarithms, and the real and imaginary parts of the principal fixed points $L_2$,
$L_e$; no relations hold among the $m$-values themselves, and short-precision probes of
$C_1(e)$ and of the hub slope were likewise null. Together with the seventeen-digit
$C_1(e)$, this is meaningful negative evidence that the phase constants live outside the
standard constant ring---consistent with their definition through the Kneser uniformization,
and consistent with the diagonal phase $\varphi_1$ landing \emph{near} but not \emph{at}
$\operatorname{Im}L_e$.

\section{A hybrid data format for base-change computation}\label{sec:format}

The results above pin down the format in which the phases become a stable computational
primitive inside a high-precision tetration engine. The design principle is that the two
tail models of Section~\ref{sec:decay} coincide throughout the tabulated range, so no open
question blocks implementation; and that every reconstructed quantity is guarded by an
identity that is a \emph{theorem}, so residuals are error meters rather than plausibility
checks.

\begin{enumerate}
\item \textbf{Phase hub.} The groupoid gives the exact representation
$\widetilde H_{b,d}=U_b^{-1}\circ U_d$ with $U_a:=\widetilde H_{c_*,a}$: one phase map per
supported base reconstructs every ordered pair, with stable error propagation since
$\inf U_a'\ge1-\delta$ is certifiable from the mode table. Direct pair tables are optional
caches, not fundamental data.
\item \textbf{Hub column.} Store $m(b)=\mu_{e,b}$ to full precision for each supported base
($N$ values), plus the low-$k$ mode heads. Reconstruct any pair mean as
$\mu_{b,d}=m(d)-m(b)+\Delta^{(2)}+\Delta^{(3)}$ (Proposition~\ref{prop:hub}); the chain closes to
$10^{-11}$--$10^{-13}$ absolute at the displayed precision (Proposition~\ref{prop:hub}),
and a direct pair run remains available beyond that.
\item \textbf{Fourier table for the working range.} For target precisions
$\lesssim150$ digits ($k\le300$): per pair a low-$k$ head ($k\le10$--$20$ complex
coefficients) plus the universal curve, $A_k\approx C_1\lvert\ln\alpha\rvert\,\Xi(k)\,
e^{\delta k}$ (three parameters per pair, $\lvert\delta\rvert\sim10^{-3}$).
\item \textbf{Certified ladder beyond.} From $\sim150$--$200$ digits upward, evaluate
pointwise by Algorithm~\ref{alg:ladder} with the certificate of
Proposition~\ref{prop:cert}: cost one $\mathrm{sexp}_d$ plus one $\slog_b$ per evaluation at
arbitrary precision. Inversion by Proposition~\ref{prop:newton}; the reverse map is never
computed separately. A $500$-digit Fourier table is possible
(Corollary~\ref{cor:yield}: $k\approx1400$, $n_{\mathrm{lim}}\approx600$--$750$, $1$--$2$ h
one-time initialization per base) but pays off only under many evaluations per pair.
\item \textbf{Instrument guard.} Enforce $k_{\max}\le2.3\,n_{\mathrm{lim}}$
(Observation~\ref{obs:floor}); choose working precision from the floor law including the
measured prefactor.
\item \textbf{Gates as proof residuals.} Mean antisymmetry
$\mu_{b,d}+\mu_{d,b}=0$ (exact by \cite[Thm.~B]{unified}; achieved $\le10^{-38}$), groupoid
residual $\Phi_{d,b}\circ\widetilde H_{b,d}+\Phi_{b,d}=0$ (achieved $2.4\times10^{-34}$,
limited only by table storage precision), and the truncation gate of
Proposition~\ref{prop:cert} (one extra tower level, $<10^{-120}$). Thresholds derive from the
certificate constants, not from experience.
\end{enumerate}

\section{Discussion and open problems}\label{sec:discussion}

The phase program set out in \cite{phaselaw} asked whether mixed-base tetration heights are
related by a small, structured, computable correction. At the level of practice the answer
is now complete: the correction is computable at arbitrary precision with certificates, its
inverse is free, its mean obeys an exact hub calculus, and its Fourier representation is
governed by one universal curve plus a handful of per-pair constants. At the level of
structure, the surprise of the campaign is how much is \emph{universal}: the decay profile,
the singularity height, strength and type are shared across every measured pair, including
pairs with no base in common; all pair dependence concentrates in one amplitude
$C_1\lvert\ln\alpha\rvert$, one abscissa $x_*(b,d)$, and one hub potential $m(b)$. We record
the open problems in the order we consider them decisive.

\begin{enumerate}
\item \textbf{The residual $\sigma_*$ decimal and the fit prefactor.} The
window-sequence verdict $\sigma_*=0.058\pm0.006$ (Section~\ref{sec:singularity}) should be
sharpened by sector-resolved apex sweeps at $\Delta x=5\times10^{-4}$, closing the seam
subtlety by evaluating the zone in the adjacent cell via $D(1+\theta)=b^{D(\theta)}$; and
the implausibly large effective exponent $\gamma\approx9.5$ (i.e.\ $\rho\approx20$ via
$\gamma=\rho/2-3/4$) quantifies that the fitted prefactor is still pre-asymptotic---the
asymptotic $(\sigma_*,c,\gamma)$ triple remains to be pinned.
\item \textbf{Prove the punctured strip.} The complexified $\sigma$-recursion
\eqref{eq:sigmarec} bounds all but finitely many levels on any substrip, so the channel can
only break at finitely many explicit levels: Conjecture~\ref{conj:universal} reduces to
controlling those, in the spirit of \cite[Rem.~7.3]{unified}, plus a local analysis at the
breach. The measured breach geometry (Observation~\ref{obs:breach}) tells us exactly where.
\item \textbf{Extend the kernel theorem.} The first-order kernel is now
characterized and constructed (Paper~V \cite{paperV}); the natural extensions are the
higher-order kernels $H_d,K_d$ through the same boundary problem (each order solves the
same cohomological left-hand side with an explicit right-hand side), certified enclosures
for the constructed values, and a self-contained verification of the classical strip
geometry (K1) of the Kneser branch consumed there.
\item \textbf{Second structure and the $x_*$ map.} Chart the indicated second singularity
level at $y\approx0.17$--$0.22$, and the pair map $x_*(b,d)$ (cheap: it is the accumulation
point of the mode phases), including the suspiciously round measured centers.
\item \textbf{Edge and large-base asymptotics of the hub.} Push the edge fit window toward
$\eta$ to test the derived coefficient $\pi\sqrt{2\eta/e}$ below $0.5\%$
(Lemma~\ref{lem:edge}), and settle $m(b)-\ln\ln b\to0$ with a rate.
\item \textbf{Foundations.} Derive the channel and peeling-growth hypotheses of
\cite{phaselaw} from the Kneser--Paulsen--Cowgill construction (the single remaining
conditional element, cf.\ \cite[Rem.~7.4]{unified}), including a stopping-rule formulation
covering skew pairs; and prove Hypothesis~\ref{hyp:basediff} from base holomorphy
\cite{paulsen}.
\end{enumerate}

Beyond these, the applied horizon that motivated the program remains: with the conversion
maps certified, a single master engine in one base plus phase tables can, in principle,
replace per-base Kneser constructions for multi-base workloads; whether the fixed-point
formulation of $F_{b,d}$ converges fast enough in practice is an empirical question the
hybrid format of Section~\ref{sec:format} is designed to answer.

\appendix

\section{Data tables}\label{app:tables}

All values from the July 2026 campaign; run metadata in Appendix~\ref{app:repro}.

\subsection*{A.1\quad Regenerated pair table (regular branch)}

Means $\mu_{b,d}$ exact to the digits shown (25+ digits available in the repository);
$A_1$, first-mode phase and $A_2/A_1$ from the $\mathrm{dps}=80$, $N=256$ runs. This
supersedes Table~1 of \cite{phaselaw}.

\begin{center}
\begin{tabular}{llll}
\toprule
$(b,d)$ & $\mu_{b,d}$ & $A_1$ & $A_2/A_1$ \\
\midrule
$(e,2)$ & $-1.1284037766289240098792179$ & $1.4368\times10^{-4}$ & $0.00968$ \\
$(2,e)$ & $+1.1284037766289240098792179$ & $1.4367\times10^{-4}$ & $0.00914$ \\
$(2,10)$ & $+2.2649608387^{\,*}$ & $4.6080\times10^{-4}$ & $0.00865$ \\
$(10,2)$ & $-2.2649608387^{\,*}$ & $4.6083\times10^{-4}$ & $0.01015$ \\
$(e,10)$ & $+1.1365570557133264408993760$ & $3.1723\times10^{-4}$ & $0.00885$ \\
$(10,e)$ & $-1.1365570557133264408993760$ & $3.1722\times10^{-4}$ & $0.00981$ \\
$(3,5)$ & $+0.5788350021355712162374735$ & --- & --- \\
\bottomrule
\end{tabular}
\end{center}

\noindent(Starred values are hub-reconstructed, $\mu=m(d)-m(b)+\Delta$; the digits shown are
limited by the two-digit published defect $\Delta(2,10)=6.4\times10^{-9}$, while the full
$\Delta^{(2)}+\Delta^{(3)}$ chain in the repository carries them to $\sim10^{-15}$. Direct-run
values (the $e$-pairs and $(3,5)$) are exact to the digits shown. The $(3,5)$ mode head
resides in the repository; amplitudes for the closer pair are proportionally smaller by the
amplitude law.)

\subsection*{A.2\quad Local decay rates, $(e,2)$, clean to $k=290$}

\begin{center}
\begin{tabular}{rl@{\qquad}rl@{\qquad}rl}
\toprule
$k$ & $\sigma_\eff$ & $k$ & $\sigma_\eff$ & $k$ & $\sigma_\eff$\\
\midrule
1 & 0.740 & 30 & 0.232 & 110 & 0.1689\\
2 & 0.567 & 40 & 0.2154 & 120 & 0.1657\\
3 & 0.487 & 50 & 0.2036 & 140 & 0.1601\\
5 & 0.404 & 60 & 0.1946 & 200 & 0.1484\\
8 & 0.344 & 70 & 0.1875 & 240 & 0.1429\\
20 & 0.260 & 80 & 0.1817 & 290 & 0.1375\\
 & & 90 & 0.1768 & & \\
 & & 100 & 0.1726 & & \\
\bottomrule
\end{tabular}
\end{center}

\noindent Saddle-model reconstruction:
$\sigma_\eff(k)=\sigma_*+1.090/\sqrt k-1.512/k$ reproduces this table to $10^{-4}$ with the
fitted parameters of Observation~\ref{obs:universal}. Digit yield: $D(60)=47$,
$D(100)=67$, $D(120)=76$, $D(300)\approx148$.

\subsection*{A.3\quad Breach zones at $y=0.085$ (complexified ladder)}

\begin{center}
\begin{tabular}{lccc}
\toprule
pair & zone $x$-interval & width & center $x_*$ (mod 1)\\
\midrule
$(e,2)$ & $[0.91,\,0.98]$ & $\approx0.07$ & $\approx0.945$\\
$(3,5)$ & $[1.01,\,1.08]$ & $\approx0.07$ & $\approx1.045$\\
$(2,10)$ & $[0.650,\,0.710]$ & $\approx0.06$ & $\approx0.680$\\
\bottomrule
\end{tabular}
\end{center}

\noindent Apex $(e,2)$: $z_*\approx(1.02\pm0.02)+(0.055\pm0.005)i$; center slope
$dx/dy\approx-2.8$; fitted singularity data $(\sigma_*,c,\gamma)$:
$(0.0789,\,59.78,\,9.47)$ for $(e,2)$ and $(0.0791,\,59.77,\,9.59)$ for $(3,5)$; phase
strength ratio $c_{\text{phase}}/c_{\text{amp}}\approx4.38$, $\chi\approx0.68$ rad
(pair-universal to $0.3\%$).

\subsection*{A.4\quad Hub raster and first-order constants}

Hub values $m(b)=\mu_{e,b}$ as in Section~\ref{sec:hub}; $24$-point raster over
$[1.46,100]$ at $40$ digits in the repository. Edge fit on $[1.46,1.6]$:
$m(b)=-3.3148\,(b-\eta)^{-1/2}-0.398\ln(b-\eta)+3.163$ (rms $0.004$);
derived coefficient $\pi\sqrt{2\eta/e}=3.23907$ (ratio $1.023$). Kernel constants:
$C_1(e)=0.00038859729244516473$; $\overline G_e=2.191463277066$;
$C_1(d)\times10^4=3.953,\,3.886,\,3.868,\,3.793,\,3.720$ for $d=2,e,3,5,10$;
diagonal ratios $A_2/A_1=0.0093355030$, $A_3/A_1=2.66505\times10^{-4}$,
$A_4/A_1=1.26473\times10^{-5}$; diagonal phase $\varphi_1=1.33536363153$.

\subsection*{A.5\quad Branch-robustness budget}

Mean shifts on the explicit $C^1$ branches versus the regular branch:
$\Delta\mu=\mu_{\mathrm{toy}}-\mu_{\mathrm{reg}}$ equals $+6.37\times10^{-4}$ for $(e,2)$,
$+7.34\times10^{-4}$ for $(3,5)$, $+3.758\times10^{-3}$ for $(e,10)$,
$+4.341\times10^{-3}$ for $(2,10)$: empirically $\lvert\Delta\mu\rvert\approx(5$--$12)\times
A_1(b,d)$---the choice of branch moves the mean by a few phase amplitudes, while the profile
shapes remain similar. This quantifies the branch sensitivity of every constant in this
paper.

\section{Reproducibility}\label{app:repro}

\emph{Pipeline.} Phases are evaluated by Algorithm~\ref{alg:ladder} against the author's
engine (fork of \texttt{fatou.gp} \cite{fatou}; PARI/GP; engine, data, scripts and run
metadata at \url{https://github.com/Lightrunnerwastaken/gp-tetration}; this paper is archived at \url{https://doi.org/10.5281/zenodo.21363593}, the companion papers at the DOIs given in the bibliography and the Authorship Statement), read-only, with per-point closure
bounds logged; modes by direct DFT on $N$-point grids; the complexified scan by the
bottom-up solver of Section~\ref{sec:singularity} with path continuation and mid-strip
cross-validation against the mode series.

\emph{Gates (all enforced per run).} (i) truncation gate: one extra explicit tower level
changes no phase value by more than $10^{-120}$; (ii) mean antisymmetry
$\le10^{-39}$ pairwise; (iii) groupoid residual $\le2.4\times10^{-34}$ (storage-limited);
(iv) toy-branch control: measured $N^{-3}$ aliasing against the independent value of
\cite[Rem.~7.5]{unified}, and inverse identity to $5.7\times10^{-101}$; (v) instrument
guard $k\le2.3\,n_{\mathrm{lim}}$.

\emph{Runs (selection).}

\begin{center}
\begin{tabular}{llrrrl}
\toprule
tag & pairs & dps & $N$ & $K$ & notes\\
\midrule
toy & $(e,2),(2,e),(2,10),(e,10),(3,5)$ & mp100 & 512 & 64 & calibration, $p\approx3$\\
reg & all 12 regular pairs & 80 & 256 & 64 & cache-hot, $<5$ s each\\
deep & $(e,2)$ & 200 & 512 & 176 & floor study, $n_{\mathrm{lim}}=30/60$\\
edge & $(1.5,e),(e,1.5)$ & 60 & 256 & 48 & init $1.2$ s\\
decisive & $(e,2)$, $n_{\mathrm{lim}}=120$ & 190 & 1024 & 300 & inits $239$ s $+$ $196$ s\\
decisive & $(3,5)$, $n_{\mathrm{lim}}=120$ & 190 & 1024 & 300 & tail universality\\
diag & $(d',e)$, $\ln d'=1+\varepsilon$ & 200 & 512 & --- & $\varepsilon=10^{-12},10^{-15}$\\
\bottomrule
\end{tabular}
\end{center}

\section*{Authorship and AI Contribution Statement}
\addcontentsline{toc}{section}{Authorship and AI Contribution Statement}

This manuscript is the product of a directed human--AI collaboration, and the division of
labor should be stated plainly. The author conceived the mixed-base phase program and its
central object $\Phi_{b,d}$, built and maintains the high-precision tetration engine (a fork
of Levenstein's \texttt{fatou.gp}, published at \url{https://github.com/Lightrunnerwastaken/gp-tetration}) against which every number in this paper was computed,
and directed the research: posing the questions, setting priorities between rounds,
auditing intermediate results, and deciding which claims met the acceptance gates. The
majority of the mathematical development and of the numerical work was produced by
artificial-intelligence systems under that direction: an interactive assistant (Claude,
Anthropic) contributed the theoretical framework of Sections~\ref{sec:ladder},
\ref{sec:dictionary}, \ref{sec:singularity}--\ref{sec:hub} (including the certificate
propositions, the saddle-point and crossover analysis, the response-kernel theory, the
third-order hub defect and the parabolic edge law), the round-by-round experimental designs,
and the review of the companion manuscripts; an autonomous research agent designed,
implemented and executed the measurement pipeline, discovered and quantified the instrument
systematics, performed the complexified scans, and produced all data reported here. The
companion manuscripts \cite{phaselaw,unified} were developed in the same directed mode with
a further AI system (ChatGPT, OpenAI) as the principal drafting contributor. In line with
prevailing journal policies, AI systems are not listed as authors; the author assumes full
and sole responsibility for the correctness of all statements, and readers should weigh the
provenance described here in their assessment. The five papers of the series are permanently archived on Zenodo under the concept DOIs \url{https://doi.org/10.5281/zenodo.21346803} (I), \url{https://doi.org/10.5281/zenodo.21361967} (II), \url{https://doi.org/10.5281/zenodo.21362206} (III), \url{https://doi.org/10.5281/zenodo.21363593} (IV), and \url{https://doi.org/10.5281/zenodo.21363809} (V). The author's role, in short, was closer to
that of a research director than of a sole working mathematician; the author considers it a
matter of scientific honesty that this be recorded rather than obscured.

\begin{thebibliography}{10}

\bibitem{phaselaw}
J.~Justus,
\emph{A phase law for mixed-base tetration},
Paper~I of the mixed-base tetration series, preprint, July 2026.
\url{https://github.com/Lightrunnerwastaken/gp-tetration}. DOI: \url{https://doi.org/10.5281/zenodo.21346803}.

\bibitem{unified}
J.~Justus,
\emph{The mixed-base tetration phase map is a circle diffeomorphism},
Paper~II of the mixed-base tetration series, preprint, July 2026.
\url{https://github.com/Lightrunnerwastaken/gp-tetration}. DOI: \url{https://doi.org/10.5281/zenodo.21361967}.

\bibitem{paperIII}
J.~Justus,
\emph{Stopped channels and offset alignment: de-conditionalizing the real-channel hypotheses for mixed-base tetration},
Paper~III of the mixed-base tetration series, preprint, July 2026.
\url{https://github.com/Lightrunnerwastaken/gp-tetration}. DOI: \url{https://doi.org/10.5281/zenodo.21362206}.

\bibitem{paperV}
J.~Justus,
\emph{The base-derivative of the Kneser family: selection principle, uniqueness, and existence via a Riemann--Hilbert problem},
Paper~V of the mixed-base tetration series, preprint, July 2026.
\url{https://github.com/Lightrunnerwastaken/gp-tetration}. DOI: \url{https://doi.org/10.5281/zenodo.21363809}.

\bibitem{kneser}
H.~Kneser,
\emph{Reelle analytische L\"osungen der Gleichung $\varphi(\varphi(x))=e^x$ und verwandter
Funktionalgleichungen},
J.\ reine angew.\ Math.\ \textbf{187} (1950), 56--67.

\bibitem{kouznetsov}
D.~Kouznetsov,
\emph{Solution of $F(z+1)=\exp(F(z))$ in complex $z$-plane},
Math.\ Comp.\ \textbf{78} (2009), 1647--1670.

\bibitem{paulsencowgill}
W.~Paulsen and S.~Cowgill,
\emph{Solving $F(z+1)=b^{F(z)}$ in the complex plane},
Adv.\ Comput.\ Math.\ \textbf{43} (2017), 1261--1282.

\bibitem{paulsen}
W.~Paulsen,
\emph{Tetration for complex bases},
Adv.\ Comput.\ Math.\ \textbf{45} (2019), 243--267.

\bibitem{fatou}
S.~Levenstein,
\emph{\texttt{fatou.gp}: high-precision Kneser sexp/slog in PARI/GP},
Tetration Forum, thread 1017, 2011--, software; used here through the fork at
\url{https://github.com/Lightrunnerwastaken/gp-tetration}.

\bibitem{olver}
F.~W.~J.~Olver,
\emph{Asymptotics and Special Functions},
Academic Press, 1974.

\bibitem{milnor}
J.~Milnor,
\emph{Dynamics in One Complex Variable},
3rd ed., Princeton University Press, 2006.

\end{thebibliography}

\end{document}
